\begin{table}[htbp]\centering\footnotesize
\caption{Perímetro de exclusión de la regla fiscal, 2018--2027}\label{cua:regla}
\begin{tabular}{llrrr}
\toprule
Ejercicio & Perímetro que propone la ILIF & Tope & Balance & Sin inversión \\
\midrule
2018 & gasto de inversión del Gobierno Federal y de las empresas … & 2,0 & $-$2,0 & 0,0 \\
2019 & gasto de inversión del Gobierno Federal y de las empresas … & 2,0 & $-$2,0 & 0,0 \\
2020 & gasto de inversión del Gobierno Federal y de las empresas … & 2,0 & $-$2,1 & $-$0,1 \\
2021 & gasto de inversión del Gobierno Federal y de las empresas … & 2,2 & $-$2,9 & $-$0,7 \\
2022 & gasto de inversión del sector público presupuestario & 3,1 & $-$3,1 & 0,0 \\
2023 & gasto de inversión del sector público presupuestario aprob… & \textbf{sin tope} & $-$3,6 & 0,0 \\
2024 & \textbf{cláusula ausente} & n.d. & $-$4,9 & $-$1,7 \\
2025 & \textbf{cláusula ausente} & n.d. & n.d. & n.d. \\
2026 & inversión física + inversión financiera + inversión en des… & 3,6 & $-$3,6 & \textbf{no publicado} \\
2027 & \textbf{cláusula ausente}: la exclusión migra a la LFIIEDB & n.d. & $-$3,4 & \textbf{no publicado} \\
\bottomrule
\end{tabular}
\\[2pt]\begin{minipage}{0.92\textwidth}\scriptsize Tope y balance en \% del PIB. Fuente: artículo 1o. de las diez iniciativas de Ley de Ingresos y los diez CGPE. La columna «sin inversión» es el balance que la regla mide de verdad; \textbf{dejó de publicarse en 2026 y sigue sin publicarse en 2027}.\end{minipage}
\end{table}
